\section{Defects and Wearout in OxRAM}
\label{sec:bg-defects}

Repeated SET and RESET cycling in OxRAM devices leads to gradual
degradation of the conductive filament, resulting in a reduced
resistance ratio $R_\text{OFF}/R_\text{ON}$~\cite{chen2011physical,
	hayakawa2017resolving}. This degradation arises from cumulative defect
generation and migration: repeated cycling creates trap states in the
switching layer, alters the local oxygen vacancy distribution, and
progressively modifies the filament morphology~\cite{chen2011physical}.
Over time, these changes increase the number of PV iterations required
to complete a transition and shift the write latency distribution toward
longer completions. The endurance of an OxRAM device is typically
specified as the number of SET/RESET cycles before the RBER exceeds a
target threshold~\cite{hayakawa2017resolving,
	fukuyama2019comprehensive}.

In addition to long-term endurance degradation, OxRAM devices exhibit a
short-term effect called filament relaxation: a transient drift in cell
resistance that occurs shortly after a RESET
operation~\cite{li2016short}. During relaxation, oxygen ions partially
diffuse back toward the filament region, causing the HRS resistance to
decrease over time before stabilizing. This effect can compromise
short-term data integrity if verification is performed too soon after
programming~\cite{li2016short}. Relaxation effects are typically
mitigated in WRITE-VERIFY schemes by introducing a delay between the
programming pulse and the resistance measurement, or by setting a
conservative verification threshold that accounts for expected
resistance drift.

Wearout has direct implications for ReRAM-based entropy sources. In the
early stages of device life, increased PV cycle counts broaden the write
latency distribution and can increase the min-entropy of individual
timing measurements~\cite{fukuyama2019comprehensive}. However, continued
degradation shifts the distribution toward longer completions, reducing
the spread of the distribution and eventually collapsing it toward a
narrow, high-latency regime with low effective entropy. From the
perspective of NIST SP~800-90B evaluation, symbol encodings used for
entropy harvesting must accommodate wearout-driven distributional drift,
with sufficient bin headroom to capture the evolving support of the
distribution without saturation~\cite{sonmez2016recommendation}. The
characterization of wearout effects on entropy yield is addressed in
Chapter~\ref{chap:trng}.